% =============================================================================
\section{Energía de un sistema de cargas puntuales}\label{sec:puntuales}
% =============================================================================

\subsection{Trabajo para trasladar una carga}

Recordemos la interpretación del potencial del Tema~1~\cite{T1}. Para trasladar una carga $q$ en un campo
electrostático sin acelerarla, un agente externo debe aplicarle en cada instante una fuerza igual y opuesta a la del
campo, $\vF_{\mathrm{ext}}=-q\vE$; decimos que el proceso es \emph{reversible} o cuasiestático, porque la carga pasa por una
sucesión de estados de equilibrio y su energía cinética es en todo momento despreciable.

\begin{lema}[Trabajo reversible y potencial]\label{lem:trabajo}
El trabajo reversible que un agente externo realiza para llevar una carga $q$ desde $A$ hasta $B$ en un campo
electrostático es
\begin{equation}\label{eq:trabajo}
  \Wext(A\to B)=q\big[\phi(B)-\phi(A)\big] ,
\end{equation}
independiente del camino. En particular, traerla desde el infinito cuesta $q\,\phi(B)$.
\end{lema}

\begin{proof}
Con $\vF_{\mathrm{ext}}=-q\vE$ y $\vE=-\grad\phi$, a lo largo de cualquier curva $\Gamma$ de $A$ a $B$,
\[
  \Wext=\int_\Gamma\vF_{\mathrm{ext}}\cdot\dl=q\int_\Gamma\grad\phi\cdot\dl=q\big[\phi(B)-\phi(A)\big],
\]
por el teorema fundamental del cálculo para integrales de línea. La segunda afirmación se sigue de $\phi(\infty)=0$.
\end{proof}

\begin{definicion}[Energía electrostática]\label{def:energia}
Se llama \emph{energía electrostática} $U_e$ de un sistema de cargas al trabajo reversible que debe realizar un agente
externo para formarlo, trayendo sus cargas desde una situación inicial en la que están en reposo, en el infinito e
infinitamente alejadas entre sí. Ese trabajo queda almacenado en el sistema como energía potencial.
\end{definicion}

\subsection{La energía de interacción}

Formemos el sistema trayendo las cargas una a una (\cref{fig:cargas}\,(a)). Para no confundir quién crea el potencial con dónde se evalúa,
llamamos
\begin{equation}\label{eq:phi-de-i}
  \phi^{(i)}(\vr)\defeq\frac{q_i}{4\pi\eo\abs{\vr-\vr_i}}
\end{equation}
al potencial creado \emph{por la carga $i$} en un punto $\vr$; en particular, $\phi^{(i)}(\vr_j)=q_i/(4\pi\eo r_{ij})$.
\begin{itemize}
  \item Traer $q_1$ no cuesta nada: aún no hay ningún campo que vencer, $W_1=0$.
  \item Traer $q_2$ hasta $\vr_2$ exige vencer el campo de $q_1$: $W_2=q_2\,\phi^{(1)}(\vr_2)=\dfrac{q_2q_1}{4\pi\eo r_{21}}$.
  \item Traer $q_3$ exige vencer los campos de $q_1$ y $q_2$: $W_3=q_3\big[\phi^{(1)}(\vr_3)+\phi^{(2)}(\vr_3)\big]$.
  \item En general, la carga $j$-ésima debe vencer el campo de las $j-1$ ya colocadas:
        \[
          W_j=q_j\sum_{i=1}^{j-1}\phi^{(i)}(\vr_j)=\sum_{i=1}^{j-1}\frac{q_jq_i}{4\pi\eo r_{ji}} .
        \]
\end{itemize}

% Figura: (a) formación de un sistema de N cargas puntuales trayéndolas una a una
% desde el infinito; (b) ejemplo de las cuatro cargas en un cuadrado.
\begin{figure}[t]
\centering
\begin{subfigure}[b]{0.5\textwidth}\centering
\begin{tikzpicture}[scale=1.0]
  % ejes
  \coordinate (O) at (0,0);
  \draw[tenue, line width=0.45pt, -{Stealth[length=1.6mm]}] (O) -- (2.9,0) node[below, etiqt] {$y$};
  \draw[tenue, line width=0.45pt, -{Stealth[length=1.6mm]}] (O) -- (0,2.4) node[left, etiqt] {$z$};
  \draw[tenue, line width=0.45pt, -{Stealth[length=1.6mm]}] (O) -- (-1.3,-1.0) node[left, etiqt] {$x$};
  % cargas ya colocadas
  \foreach \p/\n in {{(1.3,1.7)}/1,{(2.3,0.9)}/2,{(0.9,-0.9)}/3} {
    \fill[cargapos] \p circle (0.075);
    \draw[campoP, line width=0.6pt, -{Stealth[length=1.6mm]}] (O) -- \p;
  }
  \node[etiq, anchor=south west] at (1.35,1.72) {$q_1$};
  \node[etiq, anchor=west] at (2.38,0.92) {$q_2$};
  \node[etiq, anchor=north west] at (0.95,-0.95) {$q_3$};
  \node[etiq, text=campoP] at (0.42,1.05) {$\vr_1$};
  \node[etiq, text=campoP] at (1.42,0.3) {$\vr_2$};
  \node[etiq, text=campoP] at (0.25,-0.62) {$\vr_3$};
  % carga N trayéndose desde el infinito
  \fill[cargapos] (-1.4,1.55) circle (0.075);
  \node[etiq, anchor=south] at (-1.4,1.62) {$q_N$};
  \draw[campoE, line width=0.6pt, dashed, -{Stealth[length=1.8mm]}] (-2.9,2.25) -- (-1.5,1.6);
  \node[etiqt, anchor=south] at (-2.55,2.3) {desde $\infty$};
  \draw[campoP, line width=0.6pt, -{Stealth[length=1.6mm]}] (O) -- (-1.33,1.48);
  \node[etiq, text=campoP] at (-0.95,0.62) {$\vr_N$};
  \node[etiqt] at (1.9,-0.25) {$\cdots$};
\end{tikzpicture}
\caption{Formación del sistema carga a carga}
\end{subfigure}\hfill
\begin{subfigure}[b]{0.44\textwidth}\centering
\begin{tikzpicture}[scale=1.25]
  \draw[tenue, line width=0.45pt, -{Stealth[length=1.6mm]}] (-0.3,0) -- (2.2,0) node[below, etiqt] {$x$};
  \draw[tenue, line width=0.45pt, -{Stealth[length=1.6mm]}] (0,-0.3) -- (0,2.0) node[left, etiqt] {$y$};
  \draw[equipot] (0,0) -- (1.5,0) -- (1.5,1.5) -- (0,1.5) -- cycle;
  \draw[acentoB, line width=0.55pt] (0,0) -- (1.5,1.5);
  \node[etiqf, text=acentoB, font=\footnotesize, rotate=45] at (0.62,0.82) {$\sqrt2\,a$};
  \foreach \p in {(0,0),(1.5,0),(1.5,1.5),(0,1.5)} \fill[cargapos] \p circle (0.07);
  \node[etiq, anchor=north east] at (-0.03,-0.03) {$q$};
  \node[etiq, anchor=north] at (1.5,-0.08) {$q$};
  \node[etiq, anchor=south west] at (1.53,1.53) {$q$};
  \node[etiq, anchor=east] at (-0.08,1.5) {$q$};
  \draw[cota] (0,-0.55) -- (1.5,-0.55) node[midway, below, etiq] {$a$};
\end{tikzpicture}
\caption{Cuatro cargas en un cuadrado}
\end{subfigure}
\caption{Energía de un sistema de cargas puntuales. (a) El sistema se forma trayendo las cargas una a una desde el
infinito, donde están en reposo e infinitamente alejadas entre sí; la carga $q_N$ debe vencer el campo de las $N-1$ ya
colocadas. (b) Cuatro cargas iguales en los vértices de un cuadrado de lado $a$ (\cref{ej:cuadrado}): cada carga
tiene dos vecinas a distancia $a$ y una a distancia $\sqrt2\,a$.}
\label{fig:cargas}
\end{figure}


\begin{proposicion}[Energía de un sistema de cargas puntuales]\label{prop:puntuales}
La energía electrostática de un sistema de $N$ cargas puntuales $q_i$ situadas en $\vr_i$ es
\begin{equation}\label{eq:puntuales}
  U_e=\sum_{j=2}^{N}\sum_{i=1}^{j-1}\frac{q_iq_j}{4\pi\eo r_{ij}}
     =\frac12\sum_{i=1}^{N}\sum_{\substack{j=1\\ j\neq i}}^{N}\frac{q_iq_j}{4\pi\eo r_{ij}}
     =\boxed{\;\frac12\sum_{i=1}^{N}q_i\,\phi_i\;},
\end{equation}
donde
\[
  \phi_i\defeq\sum_{j\neq i}\phi^{(j)}(\vr_i)=\sum_{j\neq i}\frac{q_j}{4\pi\eo r_{ij}}
\]
es el potencial \emph{en el punto} $\vr_i$ creado por todas las cargas salvo la propia. El resultado
no depende del orden en que se traigan las cargas.
\end{proposicion}

\begin{proof}
La primera igualdad es la suma $U_e=\sum_jW_j$ de los trabajos anteriores. Su sumando es simétrico en $i$ y $j$, de modo que
la suma sobre los pares con $i<j$ es la mitad de la suma sobre todos los pares ordenados con $i\neq j$, que cuenta cada par dos
veces; esto da la segunda igualdad, en la que ya no aparece ningún orden. Finalmente, agrupando los términos según la carga $q_i$ que recibe el potencial,
\[
  U_e=\frac12\sum_{i=1}^Nq_i\sum_{j\neq i}\frac{q_j}{4\pi\eo r_{ij}}=\frac12\sum_{i=1}^Nq_i\sum_{j\neq i}\phi^{(j)}(\vr_i)=\frac12\sum_{i=1}^Nq_i\,\phi_i . \qedhere
\]
\end{proof}

\begin{observacion}[Una energía de interacción]\label{obs:interaccion}
\begin{enumerate}
  \item El factor $\frac12$ no es un descuido: en $\sum_iq_i\phi_i$ cada pareja aparece dos veces, una como $q_i\,\phi^{(j)}(\vr_i)$ y otra como
        $q_j\,\phi^{(i)}(\vr_j)$, que valen lo mismo. Con él, $\phi_i$ incluye el potencial de \emph{todas} las demás cargas, no solo el de las que ya estaban
        presentes al traer la $i$-ésima.
  \item Se excluyen los términos $i=j$: una carga no ejerce fuerza sobre sí misma. La expresión \eqref{eq:puntuales} no incluye,
        por tanto, la energía necesaria para formar cada carga puntual; las cargas «ya existían» en el infinito. Por eso se llama
        \emph{energía de interacción}.
  \item Puede ser positiva o negativa. Para dos cargas, $U_e=q_1q_2/(4\pi\eo r_{12})$ es positiva si se repelen (hubo que empujarlas
        para acercarlas) y negativa si se atraen (el campo hizo el trabajo y el sistema cedió energía).
\end{enumerate}
\end{observacion}

\begin{ejemplo}[Cuatro cargas en un cuadrado]\label{ej:cuadrado}
Sean cuatro cargas $q$ en los vértices de un cuadrado de lado $a$ (\cref{fig:cargas}\,(b)).

\emph{Primer método: trayéndolas por orden.} La primera no cuesta nada; la segunda, colocada a distancia $a$ de la primera,
$q\cdot kq/a$; la tercera, en el vértice opuesto a la primera, está a distancia $a$ de la segunda y $\sqrt2\,a$ de la primera; la
cuarta está a distancia $a$ de dos de ellas y $\sqrt2\,a$ de la otra:
\[
  U_e=0+\frac{kq^2}{a}+kq^2\Big(\frac1a+\frac{1}{\sqrt2\,a}\Big)+kq^2\Big(\frac2a+\frac{1}{\sqrt2\,a}\Big)
     =\frac{q^2}{4\pi\eo a}\Big(4+\frac{2}{\sqrt2}\Big)=\frac{q^2}{4\pi\eo a}\big(4+\sqrt2\big).
\]
\emph{Segundo método: con la fórmula simétrica.} Por simetría, las cuatro cargas ven el mismo potencial, creado por dos vecinas a
distancia $a$ y una a distancia $\sqrt2\,a$, de modo que
\[
  U_e=\frac12\cdot4\,q\,\phi_1=2q\Big(\frac{2kq}{a}+\frac{kq}{\sqrt2\,a}\Big)=\frac{q^2}{4\pi\eo a}\big(4+\sqrt2\big).
\]
El segundo método es mucho más rápido cuando hay simetría: basta un solo potencial.
\end{ejemplo}
